\chapter*{Abstract} 
\addcontentsline{toc}{chapter}{Abstract}

Tandem mass spectrometry (MS/MS) enables high-throughput identification of proteins, yet identifying novel peptides without reference databases---\emph{de novo} peptide sequencing---remains computationally challenging. Existing state-of-the-art approaches rely on autoregressive transformers that predict amino acid sequences token by token in $\mathcal{O}(L)$ sequential steps. While accurate, autoregressive decoding suffers from error propagation, directional fragmentation asymmetry between complementary $b$- and $y$-ions, and severe latency bottlenecks. Continuous normalizing flows alleviate sequential decoding but suffer from severe discretization error upon rounding continuous latent states to discrete amino acid masses. 

In this work, we introduce \textbf{DFlowNovo}, a non-autoregressive generative framework that models \emph{de novo} peptide sequencing as continuous-time Discrete Flow Matching (DFM) over the probability simplex of amino acids, coupled with exact Dynamic Programming Knapsack reachability filtering (\texttt{KnapsackDP}). DFlowNovo generates complete peptide sequences in a constant budget of $K=20$ reverse flow steps, achieving an inference throughput of 227.1--255.8 spectra per second on an NVIDIA GPU (a 4.3$\times$ to 4.9$\times$ speedup over autoregressive beam search). On the standardized 104k Nine-Species test benchmark, DFlowNovo achieves 68.28\% strict exact match and 83.01\% amino acid F1 score, outperforming InstaNovo v1.2 (65.48\%) and accepting 86,412 peptide-spectrum matches at 80\% precision. Furthermore, we dissect the physical foundations of mass spectrometry, demonstrating that remaining prediction discrepancies on human data are predominantly driven by isomeric Leucine/Isoleucine equivalence ($113.084\,\text{Da}$) which requires radical-driven $w$-ion fragmentation to break.

\bigskip
\begin{center}
\textbf{R\'esum\'e}
\end{center}

La spectrom\'etrie de masse en tandem (MS/MS) permet l'identification des prot\'eines \`a haut d\'ebit, mais le s\'equen\c{c}age \emph{de novo}---l'identification des peptides sans base de donn\'ees de r\'ef\'erence---reste un d\'efi informatique majeur. Les mod\`eles autor\'egressifs actuels pr\'edisent les acides amin\'es de mani\`ere s\'equentielle en $\mathcal{O}(L)$ \'etapes, ce qui induit une propagation des erreurs, une asym\'etrie entre ions compl\'ementaires $b$ et $y$, et une forte latence. Les mod\`eles \`a flux continus \'evitent ce co\^ut s\'equentiel mais souffrent d'une perte de discr\'etisation critique lors de la projection sur les masses discr\`etes des r\'esidus.

Ce travail pr\'esente \textbf{DFlowNovo}, un cadre g\'en\'eratif non-autor\'egressif formulant le s\'equen\c{c}age \emph{de novo} comme un appariement de flux discret (DFM) en temps continu sur le simplexe de probabilit\'e, guid\'e par une programmation dynamique exacte du sac \`a dos (\texttt{KnapsackDP}). DFlowNovo d\'ecode les s\'equences peptidiques enti\`eres en un nombre fixe de 20 pas d'int\'egration, atteignant une vitesse de 227 \`a 256 spectres par seconde (4,3 \`a 4,9 fois plus rapide que les mod\`eles autor\'egressifs). Sur le banc d'essai biologique des Neuf Esp\`eces (104\,163 spectres), DFlowNovo atteint 68,28\% d'exactitude stricte et 83,01\% de score F1, surpassant InstaNovo v1.2 tout en identifiant 86\,412 spectres \`a 80\% de pr\'ecision. Enfin, une analyse biophysique d\'emontre que les ambigu\"it\'es r\'esiduelles d\'ecoulent de l'isom\'erie fondamentale Leucine/Isoleucine, indissociable par fragmentation collisionnelle standard.

\bigskip
\noindent\textbf{Code and Model Availability:} Open-source codebase, checkpoints, and demonstration apps are accessible at:
\begin{itemize}\setlength{\itemsep}{1pt}
\item Repository: \url{https://github.com/joelinator/JoelResearch}
\item Model Checkpoint: \url{https://huggingface.co/joelinator/dflow-novo-model}
\item Interactive Studio: \url{https://huggingface.co/spaces/joelinator/dflow-novo}
\end{itemize}

\vfill
\section*{Declaration}
I, the undersigned, hereby declare that the work contained in this research project is my original work, and that any work done by others or by myself previously has been acknowledged and referenced accordingly.

\vspace{0.3cm}
\includegraphics[height=1.8cm]{images/signature.png} \newline \hrule
\vspace{0.2cm}
Jo\"el G\'ed\'eon, September 2026
